% from ega4-2-fr.tex
\oldpage[IV-2]{77}

finie $k_1$ de $k$ telle que
$(X\otimes_k k_1)_{\mathrm{red}}$ soit séparable sur $k_1$ (4.6.6); donc
$(X\otimes_k k_1)_{\mathrm{red}}\otimes_{k_1}\overline k$ est réduit
(4.6.4). En vertu de (4.7.3.1), on est donc de nouveau ramené au cas où $k'$
est une extension finie radicielle de $k$, et on conclut comme plus haut.

\item[(iii)] Notons que deux extensions $k'$, $k''$ de $k$ sont contenues
dans une même extension algébriquement close de $k$; pour montrer que
l'ensemble des nombres $\operatorname{long}(\mathcal O_{x'_\alpha})$, pour
$k'$ parfaite, ne dépend pas du choix de $k'$, on peut donc se borner à
comparer ces nombres pour $k'$ et pour une extension algébriquement close
$k''$ de $k'$; dans ce cas, l'assertion résulte de (i), puisque $k''$ est
séparable sur $k'$. Pour montrer que tous les nombres
$\operatorname{long}(\mathcal O_{x'_\alpha})$ sont égaux, on peut donc se
borner au cas où $k'=\overline k$, clôture algébrique de $k$. Or, l'extension
radicielle finie $k_1$ de $k$ étant alors déterminée comme dans (ii), les
$x'_\alpha$ sont tous au-dessus de l'unique point générique de
$X_1=X\otimes_k k_1$, et la relation
$\operatorname{long}(\mathcal O_{x'_\alpha})=
\operatorname{long}(\mathcal O_{x_1})$ pour tout $\alpha$ résulte encore de
(4.7.3.1), vu le choix de $k_1$.
\end{enumerate}
\end{proof}

\begin{definition}[4.7.4]
\label{IV.2.4.7.4-fr}
Soient $k$ un corps, $X$ un $k$-préschéma intègre, $K$ le corps des fonctions
rationnelles sur $X$. On appelle \emph{multiplicité radicielle} de $K$ (ou de
$X$) sur $k$ la borne supérieure des longueurs des anneaux artiniens
$K\otimes_k k_1$, lorsque $k_1$ parcourt l'ensemble des extensions
radicielles finies de $k$. On appelle \emph{multiplicité séparable} de $K$ (ou
de $X$) sur $k$ le nombre géométrique des composantes irréductibles de $X$ si
ce nombre est fini, et $+\infty$ dans le cas contraire. Enfin, on appelle
\emph{multiplicité totale} de $K$ (ou de $X$) sur $k$ le produit de la
multiplicité radicielle et de la multiplicité séparable de $K$ sur $k$.
\end{definition}

On notera que si $p$ est l'exposant caractéristique de $k$ et si la
multiplicité radicielle de $K$ sur $k$ est finie, c'est une puissance $q=p^f$
de $p$, comme il résulte de (4.7.3, (ii)); lorsque $p\ne1$, $f$ est bien
déterminé et est appelé l'\emph{exposant d'inséparabilité} de $K$ sur $k$;
lorsque $p=1$, la multiplicité radicielle est toujours égale à $1$.

Dire que la multiplicité radicielle (resp. séparable, resp. totale) de $X$ sur
$k$ est $1$ signifie que $X$ est géométriquement réduit (resp.
géométriquement irréductible, resp. géométriquement intègre) sur $k$.

Lorsque $X$ est localement de type fini, la multiplicité radicielle de $X$ sur
$k$ est encore la longueur commune des anneaux locaux $A_i$ aux idéaux
premiers minimaux de l'anneau total des fractions $A$ de $K\otimes_k k'$,
pour toute extension parfaite $k'$ de $k$; la multiplicité séparable de $X$
sur $k$ est le nombre d'idéaux premiers minimaux de $A=K\otimes_k\Omega$ pour
toute extension algébriquement close $\Omega$ de $k$, et enfin la multiplicité
totale est la somme des longueurs des anneaux locaux $A_i$ de
$A=K\otimes_k\Omega$ en ses idéaux premiers minimaux; ces nombres sont tous
trois finis dans ce cas.

\begin{definition}[4.7.5]
\label{IV.2.4.7.5-fr}
Soient $k$ un corps, $X$ un $k$-préschéma localement de type fini sur $k$,
$\mathscr F$ un $\mathcal O_X$-Module cohérent, $x$ un point de $X$ tel que
$\mathscr F_x$ soit un $\mathcal O_x$-module de longueur finie. On appelle
\emph{longueur géométrique} de $\mathscr F$ en $x$ (relative à $k$), ou
\emph{multiplicité radicielle de $x$ pour $\mathscr F$} le produit
\[
\lambda_x(\mathscr F)
=\operatorname{long}_{\mathcal O_x}(\mathscr F_x)\mu_i(k(x)\mid k),
\]
où $\mu_i(k(x)\mid k)$ désigne la multiplicité radicielle de $k(x)$ sur $k$.
\end{definition}

Lorsque $X$ est intègre, et $x$ le point générique de $X$, on a
$\operatorname{long}(\mathcal O_x)=1$, donc la multiplicité radicielle de $x$
pour $\mathcal O_X$ n'est autre que la multiplicité radicielle de $X$ sur $k$,
définie dans (4.7.4).

